\usepackage{graphicx}


% For theorems and such
\usepackage{amsmath}
\usepackage{amssymb}
\usepackage{mathtools}
\usepackage{amsthm}

% \usepackage{tikz,lipsum,lmodern}
\usepackage[most]{tcolorbox}  % takeaway box
\tcbuselibrary{skins}


\usepackage{booktabs}
\usepackage{caption}
\usepackage{makecell}

% %%%%%%%%%%%%%%%%%%%%%%%%%%%%%%%%
% % THEOREMS
% %%%%%%%%%%%%%%%%%%%%%%%%%%%%%%%%
% \theoremstyle{plain}
% \newtheorem{theorem}{Theorem}[section]
% \newtheorem{proposition}[theorem]{Proposition}
% \newtheorem{lemma}[theorem]{Lemma}
% \newtheorem{corollary}[theorem]{Corollary}
% \theoremstyle{definition}
% \newtheorem{definition}[theorem]{Definition}
% \newtheorem{assumption}[theorem]{Assumption}
% \theoremstyle{remark}
% \newtheorem{remark}[theorem]{Remark}
% \newtheorem*{remark*}{Remark}
% \newtheorem{example}[theorem]{Example}

\usepackage{multirow}
\usepackage[dvipsnames]{xcolor}
\usepackage[table]{xcolor}
\usepackage[normalem]{ulem}
\usepackage{pifont}
\useunder{\uline}{\ul}{}

\usepackage{wrapfig}
\usepackage{lipsum}
\usepackage{array}
\usepackage{tabulary}
\usepackage{tabularx}
\usepackage{listings}
\usepackage{enumitem}
\usepackage{algorithm}
\usepackage{algorithmic}

\definecolor{mygreen}{HTML}{6AA66E}
\definecolor{myblue}{HTML}{5471AB}
\definecolor{myred}{HTML}{B65655}
\definecolor{myorange}{HTML}{D1885C}
\definecolor{mypurple}{HTML}{7F73AF}
\definecolor{mybrown}{HTML}{8F7963}

% 定义上标样式的命令
% \up
\newcommand{\up}[1]{$^{\textcolor{RoyalBlue}{+#1}}$} 
% \down
\newcommand{\down}[1]{$^{\textcolor{gray}{-#1}}$}

\lstdefinestyle{mystyle}{
    basicstyle=\ttfamily\footnotesize,
    breaklines=true, % Automatically break lines in the listings
    columns=flexible, % Adjust the column width for better formatting
    captionpos=b,
}
\lstset{style=mystyle}

\usepackage{subcaption}

\newcommand{\norm}[1]{\left\lVert#1\right\rVert}
\newcommand{\tabincell}[2]{\begin{tabular}{@{}#1@{}}#2\end{tabular}} 
\newcommand{\ie}{\emph{i.e., }}
\newcommand{\eg}{\emph{e.g., }}
\newcommand{\etal}{\emph{et al.}}
\newcommand{\st}{\emph{s.t. }}
\newcommand{\etc}{\emph{etc.}}
\newcommand{\wrt}{\emph{w.r.t. }}
\newcommand{\cf}{\emph{cf. }}
\newcommand{\aka}{\emph{aka. }}
% \newcommand{\vs}{\emph{v.s. }}
\newcommand{\cmark}{\ding{51}}%
\newcommand{\xmark}{\ding{55}}%
\newcommand{\boldcmark}{\ding{52}}%
\newcommand{\boldxmark}{\ding{56}}%

\newcommand{\kxnote}[1]{\textcolor{orange}{#1}}

\newcommand{\updated}[1]{\textcolor{Blue}{#1}}


\newcommand{\methodname}{method name}
\newcommand{\MethodName}{Method Name}

% %%% 调间距用的 %%%
% \setlength{\abovecaptionskip}{5pt}
% \setlength{\belowcaptionskip}{0pt}
% \setlength{\textfloatsep}{5pt}

% \newlength{\origtextfloatsep}
% \setlength{\origtextfloatsep}{\textfloatsep}

% 定义名为 'paperbox' 的新环境
% [1] 表示有1个参数
% [Summary] 表示参数的默认值是 Summary
\newtcolorbox{paperbox}[1][Summary]{
    enhanced,
    colback=blue!3,
    colframe=NavyBlue,
    % 标题挂载设置
    attach boxed title to top left={xshift=2mm, yshift=-3mm},
    boxed title style={
        colback=NavyBlue,
        rounded corners, % 如果想要完全圆角去掉这行
        arc=3pt,
        top=0pt, bottom=0pt, left=2pt, right=2pt
    },
    % 字体与标题内容
    fonttitle=\bfseries\small, % 保持和正文一致的粗体
    title={#1}, % #1 代表传入的标题参数
    % 内部边距
    top=10pt,
    bottom=6pt,
    left=4pt,
    right=4pt,
    boxsep=1pt,
    % 边框样式
    arc=3pt,
    boxrule=0.8pt,
    % 外部间距
    before skip=5pt,
    after skip=5pt,
    % 如需要阴影可注释下面这行
    drop fuzzy shadow
}